\chapter{Additional Data}
\label{app:data}

\section{Drag-Type Paddle: Power Stroke and Recovery}
\label{app:drag-split}

\begin{table}[H]
  \centering
  \caption[Power-stroke and recovery contributions]{Contributions of power stroke and recovery to the mean thrust of the fan
    paddle, \eqref{eq:split}, in \si{\milli\newton}, for $\tau_o=0$ and $\tau_c=0.33$ (V3trap) or $\tau_c=\mathrm{DUTY}$ (robot gait).
    Net: two-state composite; corrected: range of the added-mass correction \eqref{eq:addedmass}.}
  \label{tab:app-split}
  \footnotesize
  \begin{tabular}{@{}lcccc|cccc@{}}
    \toprule
    & \multicolumn{4}{c|}{V3trap} & \multicolumn{4}{c}{Robot gait} \\
    $\U$ (\si{\metre\per\second}) & power & recovery & net & corrected & power & recovery & net & corrected \\
    \midrule
    0     & 88.0 & $-5.3$  & 82.7 & 72.9--82.7 & 29.4 & $-6.0$  & 23.4    & 23.4 \\
    0.08  & 85.7 & $-10.2$ & 75.6 & 65.8--71.4 & 18.1 & $-12.9$ & 5.2     & 3.9--5.2 \\
    0.15  & 75.4 & $-17.2$ & 58.2 & 48.4--50.4 & 12.6 & $-25.7$ & $-13.1$ & $-15.5$--$-13.1$ \\
    0.223 & 62.5 & $-26.9$ & 35.7 & 23.9--25.9 & --   & --      & --      & -- \\
    0.30  & 50.9 & $-39.1$ & 11.7 & $-4.1$--1.9 & --  & --      & --      & -- \\
    \bottomrule
  \end{tabular}
\end{table}

\section{Fluke Cases}
\label{app:fluke}

\begin{table}[H]
  \centering
  \caption[All fluke cases]{All single-leg fluke cases. $\alpha_{75}$: 75th percentile of $|\alpha|$ over the cycle.}
  \label{tab:app-fluke}
  \footnotesize
  \setlength{\tabcolsep}{5pt}
  \begin{tabular}{@{}llccccc@{}}
    \toprule
    Case & Motion & $\U$ (\si{\metre\per\second}) & $\Tm$ (mN) & $\Pm$ (mW) & $\etaF$ & $\alpha_{75}$ \\
    \midrule
    L2U000 & L0 & 0     & 45.5 & 18.1 & --    & \SI{73}{\degree} \\
    L2U008 & L0 & 0.08  & 36.5 & 17.1 & 0.171 & -- \\
    L2U015 & L0 & 0.15  & 29.1 & 16.6 & 0.263 & -- \\
    L0     & L0 & 0.223 & 21.4 & 16.3 & 0.293 & \SI{19.9}{\degree} \\
    L2U030 & L0 & 0.30  & 12.7 & 15.6 & 0.245 & \SI{13.7}{\degree} \\
    L1a10  & $\alpha$ law \SI{10}{\degree} & 0.223 & $\approx31$ & $\approx35$ & $\approx0.20$ & \SI{13}{\degree} \\
    L1a28  & $\alpha$ law \SI{28}{\degree} & 0.223 & 26.2 & 23.9 & 0.244 & \SI{30}{\degree} \\
    L1a20  & $\alpha$ law \SI{20}{\degree}, paddle-leg hinge path & 0.40 & 15.4 & 24.3 & 0.25 & \SI{21.7}{\degree} \\
    L3U030 & scheduled, $\theta_0=\SI{21.6}{\degree}$ & 0.30 & 13.9 & 10.5 & 0.398 & \SI{21.0}{\degree} \\
    L3U040 & scheduled, $\theta_0=\SI{17.6}{\degree}$ & 0.40 & 10.5 & 11.4 & 0.368 & \SI{16.4}{\degree} \\
    L0c    & L0, coarse mesh & 0.223 & \multicolumn{4}{l}{$-7.3$\,\% thrust in the common window (\cref{sec:ver-mesh})} \\
    \bottomrule
  \end{tabular}
\end{table}

\section{Actuator Requirements}
\label{app:actuators}

\begin{table}[H]
  \centering
  \caption[Actuator requirements]{Peak joint speeds and peak hydrodynamic moments of the motions (propulsor only; inertia and drag
    of the linkage not included), compared with the servos of the robot.}
  \label{tab:actuators}
  \small
  \begin{tabularx}{\textwidth}{@{}l>{\raggedright\arraybackslash}X>{\raggedright\arraybackslash}X>{\raggedright\arraybackslash}X@{}}
    \toprule
    & Paddle, V3trap & Fluke, L0 & Fluke, L3 at \SI{0.30}{\metre\per\second} \\
    \midrule
    Leg joints & hip \SI{214}{\degree\per\second}, curl \SI{459}{\degree\per\second} & knee $\approx$\SI{310}{\degree\per\second} & as L0 \\
    Propulsor joint & fold within \SIrange{45}{60}{\milli\second} at $\tau_c$ & pitch \SI{712}{\degree\per\second} & pitch \SI{366}{\degree\per\second} \\
    Hydrodynamic moment & \SI{35.6}{\milli\newton\metre} about the hip & \SI{12.6}{\milli\newton\metre} about the pin (at \SI{0.30}{\metre\per\second}) & \SI{7.4}{\milli\newton\metre} about the pin \\
    \midrule
    Available & \multicolumn{3}{>{\raggedright\arraybackslash}p{0.70\textwidth}}{PTK\,7465W at \SIrange{6.0}{8.4}{\volt}: stall torque
      \SIrange{0.44}{0.57}{\newton\metre}, no-load speed \SIrange{650}{860}{\degree\per\second}; no controllable fold on the hardware} \\
    \bottomrule
  \end{tabularx}
\end{table}

\chapter{Inter-Limb Coordination in the Reduced-Order Model}
\label{app:coord}

This appendix reports the results of the reduced-order model on the coordination of the four legs. They do not affect the
comparison of drag- and lift-based propulsion in \cref{ch:mujoco}, but they show how strongly the attitude limits shape the
optimal gait.

\section{Gait Families without Lift}

With drag only and moderate limits, all five families find feasible gaits
(\cref{tab:mj-families}, \cref{fig:mj-families}a). Bound is the fastest,
$0.481\pm\SI{0.034}{\metre\per\second}$ or about two body lengths per second, with a margin larger than the seed
scatter. It is left--right symmetric, so heading and roll vanish and all degrees of freedom serve propulsion;
all three bound solutions sit at the \SI{15}{\degree} pitch limit and at the upper frequency bound of
\SI{2.5}{\hertz}. About two thirds of all solutions end within \SI{1}{\degree} of a limit, and 14 of 15 keep
their servos at the torque limit throughout. The limits shape the answer and are reported with it.

With strict limits and drag only, the advantage of bound disappears: it drops to
$0.159\pm\SI{0.020}{\metre\per\second}$, all ten family solutions sit at the \SI{5}{\degree} pitch limit, and walk,
pace and trot differ by less than the seed scatter (\cref{fig:mj-families}b). Pitch is the bottleneck of every
coordination.

\begin{table}[t]
  \centering
  \caption[Fastest gait per family]{Fastest gait per family (speed in \si{\metre\per\second}, mean $\pm$
    standard deviation over seeds). Moderate limits $10/10/\SI{15}{\degree}$, strict limits
    $10/5/\SI{5}{\degree}$. Seeds per family: 3 (drag only, moderate), 2 (drag only, strict), 5 for bound, pace and
    trot and 2 for walk and pronk (drag + lift, strict).}
  \label{tab:mj-families}
  \small
  \begin{tabular}{@{}lccc@{}}
    \toprule
    Family & Drag only, moderate & Drag only, strict & Drag + lift, strict \\
    \midrule
    bound & $\mathbf{0.481\pm0.034}$ & $0.159\pm0.020$ & $\mathbf{0.295\pm0.139}$ \\
    pace  & $0.345\pm0.064$ & $0.191\pm0.062$ & $0.204\pm0.085$ \\
    trot  & $0.344\pm0.048$ & $0.172\pm0.013$ & $0.139\pm0.051$ \\
    walk  & $0.344\pm0.098$ & $0.209\pm0.089$ & $0.143\pm0.077$ \\
    pronk & $0.265\pm0.038$ & $0.145\pm0.043$ & $0.118\pm0.064$ \\
    \midrule
    free phase, $10^4$ evaluations & -- & -- & $0.303\pm0.078$ \\
    \bottomrule
  \end{tabular}
\end{table}

\begin{figure}[t]
  \centering
  \includegraphics[width=\textwidth]{fig_mj_families}
  \caption[Fastest gaits per family in MuJoCo]{Fastest gait per family in the reduced-order model; bars are
    seed means, dots individual seeds. (a) Drag only, moderate limits. (b) Drag only, strict limits. (c) Drag and
    lift, strict limits, including three free-phase searches with $10^4$ evaluations. Free-phase searches with
    2500 evaluations are not shown because they had not converged.}
  \label{fig:mj-families}
\end{figure}

\section{Coordination with Lift under Strict Limits}

With lift and strict limits, all fifteen trot, pace and bound solutions over five seeds are lift-propelled,
feasible and do not surface (\cref{tab:mj-families}, \cref{fig:mj-families}c). Bound is again the fastest, up to
\SI{0.483}{\metre\per\second} for one seed, as fast as the fastest drag-only gait under moderate limits, but its
servos are saturated almost all the time. Trot has the lowest mean power ($1.23\pm\SI{0.90}{\watt}$) but is also the slowest of the three, so its low energy per metre mixes \enquote{economical} with \enquote{slow}.

To separate the two, the least power at equal speed was computed per family with lift and strict limits
(five seeds, \cref{fig:mj-coord}). At \SI{0.10}{\metre\per\second} trot needs \SI{0.10}{\watt} against
\SI{0.17}{\watt} for bound, and at \SI{0.15}{\metre\per\second} \SI{0.48}{} against \SI{0.51}{\watt}; the medians
over seeds are close. At \SI{0.20}{\metre\per\second} trot needs \SI{1.68}{\watt}, while pronk and bound need
\SI{0.96}{} and \SI{0.98}{\watt}, and only three of five trot seeds are feasible; trot's speed limit under strict
limits is about \SI{0.22}{\metre\per\second}. All 55 feasible solutions of this experiment are lift-propelled.
Diagonal coordination is therefore at most marginally more economical at low speed; above about
\SI{0.15}{\metre\per\second}, fore--aft coordinations are both more economical and faster.
\Cref{fig:mj-trot} shows one lift-propelled trot at \SI{0.221}{\metre\per\second} and \SI{1.03}{\watt}, whose
servos never saturate.

The free-phase search needs a larger budget. With 2500 evaluations it reached only
\SIrange{0.05}{0.07}{\metre\per\second}; with $10^4$ evaluations, three seeds reached
$0.303\pm\SI{0.078}{\metre\per\second}$, all lift-propelled and all converged to fore--aft coordinations close to
bound or pronk (deviation 15--19\,\% of a period), not to trot. The best free gait, \SI{0.369}{\metre\per\second}
at \SI{3.3}{\watt}, is faster than three of the five bound optima and needs less energy per metre than all five, which shows that better
coordinations exist outside the fixed families.

\begin{figure}[t]
  \centering
  \includegraphics[width=\textwidth]{fig_mj_trot_frames.png}\\[1mm]
  {\small (a)}\\[2mm]
  \includegraphics[width=\textwidth]{fig_mj_trot.png}\\
  {\small (b)}
  \caption[A lift-propelled trot in MuJoCo]{A lift-propelled trot found with strict limits (\SI{0.90}{\hertz},
    $\mathrm{DUTY}=0.59$, \SI{0.221}{\metre\per\second}, \SI{1.03}{\watt}). (a) Four frames of the MuJoCo animation,
    \SIrange{0.25}{0.30}{\second} apart, covering one stroke period of \SI{1.11}{\second}; the camera follows the trunk. (b) Top:
    power strokes of the four legs over the last three cycles; diagonal legs act together. Bottom: forward speed of the trunk
    after the ramp-in, mean dotted.}
  \label{fig:mj-trot}
\end{figure}


\begin{figure}[t]
  \centering
  \includegraphics[width=0.58\textwidth]{fig_mj_coord}
  \caption[Least power per gait family]{Least mechanical power to reach a required speed per gait family, with lift and strict limits
    (best of five seeds).}
  \label{fig:mj-coord}
\end{figure}

\section{Electrical Power}

Because most of the fastest gaits keep their servos at the torque limit, the mechanical-power ranking could
hide large copper losses. The electrical power was therefore recomputed for all 222 valid optima. Averaged per family, it lies within about $\pm30$\,\% of the mechanical power, because at saturation the back-EMF term cancels much of the stall torque. The exceptions are trot and pronk with drag only under moderate limits, whose electrical power is 71\,\% and 44\,\% higher (for trot \SI{19.6}{} instead of \SI{11.5}{\watt}). This moves trot behind pace and walk in energy per metre; apart from this, the statements of this appendix hold for the electrical power as well.

\chapter{Real-Time Solver: Tow Runs and Whole-Robot Trends}
\label{app:flare-trends}

This appendix collects the Flare Live results. The tow runs are the basis of the cross-check in \cref{sec:ver-flare}; the
self-propelled runs show trends of the whole robot but no reliable magnitudes.

\section{Tow Runs}
\label{app:flare}

\begin{table}[H]
  \centering
  \caption[Flare Live tow runs]{Flare Live tow runs with the L0 motion: streamwise fluid force on the fluke (fluke frame) and net
    thrust of the whole flapping leg, in \si{\milli\newton}, averaged over $t=\SIrange{3}{9}{\second}$.}
  \label{tab:app-flare}
  \footnotesize
  \begin{tabular}{@{}lcccccc@{}}
    \toprule
    $\U$ (\si{\metre\per\second}) & Lattice & Fluke, flapping & Fluke, idle & Leg net thrust & CFD fluke & $\alpha_{75}$ \\
    \midrule
    0     & \SI{3}{\milli\metre} & 20.0  & --      & 25.7 & 45.5 & \SI{73.8}{\degree} \\
    0.1   & \SI{3}{\milli\metre} & 8.5   & --      & 15.9 & $\approx34$ & \SI{41.1}{\degree} \\
    0.223 & \SI{3}{\milli\metre} & $-0.1$ & $-12.6$ & 9.6  & 21.4 & \SI{19.7}{\degree} \\
    0.223 & \SI{2}{\milli\metre} & 4.1   & $-11.6$ & 5.8  & 21.4 & \SI{19.8}{\degree} \\
    \bottomrule
  \end{tabular}
\end{table}

\section{Self-Propelled Runs}

Because of the thrust deficit of \cref{sec:ver-flare}, the self-propelled Flare Live runs are used only for trends
(\cref{fig:flare-robot}, \cref{tab:flare}). With the resistance $K=\SI{0.65}{\newton\second\squared\per\metre\squared}$
($D_\text{mid}$), a trot at \SI{2}{\hertz} with a target angle of attack of \SI{18}{\degree} swims at
\SI{0.154}{\metre\per\second}, about half the CFD-based prediction, consistent with the missing fluke thrust.
\begin{itemize}
  \item \emph{Frequency.} Speed grows with frequency as $\U\propto f^{0.71}$ at $K=0.65$ ($f^{0.51}$ at $K=3.95$), while power grows
        as $f^{1.8}$; the cost of transport rises with frequency. The Strouhal number stays nearly constant.
  \item \emph{Angle of attack.} The target \SI{18}{\degree} gives the highest speed and the lowest cost of transport; \SI{28}{\degree}
        is slightly slower, \SI{10}{\degree} much slower. The controller does hold the target: in the \SI{10}{\degree} run, the
        mid-stroke angle of attack has a median of \SI{10.1}{\degree}, but the fluke has to rotate by \SI{122}{\degree} per cycle and
        sits at its end stop 11\,\% of the time. In CFD, the \SI{10}{\degree} law gave the most thrust at this speed
        (\cref{sec:res-lift}). The penalty for small angles in Flare Live is consistent with its unresolved lift; only the statement
        that efficiency is highest near \SIrange{18}{20}{\degree} is supported by both solvers.
  \item \emph{Resistance.} Lowering $K$ from 0.65 to 0.27 (bare hull) raises the speed only from \SI{0.154}{} to
        \SI{0.160}{\metre\per\second}, so in Flare Live the robot is limited by its thrust rather than by its hull.
  \item \emph{Attitude.} The hull heaves by less than \SI{0.6}{\milli\metre} and pitches by less than \SI{0.3}{\degree} in the trot:
        diagonal coordination keeps the floating robot level.
\end{itemize}

\begin{table}[t]
  \centering
  \caption[Flare Live self-propelled runs]{Self-propelled Flare Live runs, trot, $K=\SI{0.65}{\newton\second\squared\per\metre\squared}$.
    COT $=P_\text{j}/(mg\U)$ with the joint-drive power $P_\text{j}$ and $m=\SI{0.863}{\kilo\gram}$.}
  \label{tab:flare}
  \small
  \begin{tabular}{@{}lccc@{}}
    \toprule
    Run & $\U$ (\si{\milli\metre\per\second}) & $P_\text{j}$ (mW) & COT \\
    \midrule
    \SI{2}{\hertz}, $\alpha$ \SI{10}{\degree} & 100.2 & 261 & 0.31 \\
    \SI{2}{\hertz}, $\alpha$ \SI{18}{\degree} & \textbf{153.8} & 239 & \textbf{0.18} \\
    \SI{2}{\hertz}, $\alpha$ \SI{28}{\degree} & 143.8 & 249 & 0.20 \\
    \SI{2.5}{\hertz}, $\alpha$ \SI{18}{\degree} & 183.1 & 355 & 0.23 \\
    \SI{3}{\hertz}, $\alpha$ \SI{18}{\degree} & 205.0 & 488 & 0.28 \\
    \bottomrule
  \end{tabular}
\end{table}

\begin{figure}[t]
  \centering
  \includegraphics[width=\textwidth]{fig_flare_robot}
  \caption[Flare Live free-swimming trends]{Self-propelled robot in Flare Live (trot). (a) Speed versus stroke frequency for two
    hull resistances. (b) Speed versus the target angle of attack at \SI{2}{\hertz}. (c) Cost of transport versus frequency.}
  \label{fig:flare-robot}
\end{figure}

The joint-drive power in Flare Live (\SIrange{240}{490}{\milli\watt}) is an order of magnitude above the hydrodynamic power of
the CFD (about \SI{65}{\milli\watt} for four flukes), because it includes the power to move the linkages, their inertia and the
controller effort.

\chapter{Reproducibility}
\label{app:repro}

\section{Software}

\begin{sloppypar}
\begin{itemize}
  \item OpenFOAM v2606, \texttt{overPimpleDyMFoam}; meshes with \texttt{blockMesh} and \texttt{snappyHexMesh}. Case generators,
        run scripts and post-processing in Python (\texttt{make\_leg\_cfd.py}, \texttt{make\_fluke\_cfd.py},
        \texttt{fluke\_traj.py}, \texttt{fluke\_forces.py}).
  \item MuJoCo 3.14 with the Python package \texttt{cma}; the gait-optimization code \texttt{swimopt} (\url{https://github.com/lxtuni/swimopt}), including the raw results of all
        234 runs, 222 of them valid (\texttt{thesis\_notes/data/all\_runs.csv}). The replay of the drag-only optima with lift
        (\cref{sec:mj-liftshare}) is documented in the notes of the repository, not in \texttt{all\_runs.csv}.
  \item Flare Live with the scripts \texttt{run\_float.py} (self-propelled runs) and \texttt{bench\_foil.py} (tow runs).
\end{itemize}
\end{sloppypar}
\TODO{Add the repository of the CFD and Flare Live scripts and the commit hashes of the final versions.}

\section{MuJoCo Settings}

\begin{table}[H]
  \centering
  \caption[MuJoCo settings]{Settings of the reduced-order gait optimization.}
  \label{tab:app-mujoco}
  \small
  \begin{tabularx}{\textwidth}{@{}l>{\raggedright\arraybackslash}X@{}}
    \toprule
    Item & Value \\
    \midrule
    Time step, integrator & \SI{2}{\milli\second}, implicit-fast \\
    Evaluation & settle \SI{1.0}{\second}, ramp-in \SI{1.5}{\second}, score \SI{8}{\second} rounded up to whole periods \\
    Servo & position servo, stiffness 60, critical damping, $\tau_s=\SI{3}{\newton\metre}$, $\omega_\text{max}=\SI{400}{\degree\per\second}$ \\
    Flipper coefficients & $C_D=(2.2,0.10,0.10)$, $C_a=1.0$, $c_v=0.05$, $c_l=1.1$ \\
    Link coefficients & $C_D=(0.5,0.5,0.25)$, $C_a=0.3$, $c_v=0.02$ \\
    Trunk coefficients & $C_D=(0.9,1.6,1.8)$, $C_a=0.2$, $c_v=0.10$ \\
    Gait & $H=2$; $f\in[0.2,2.5]$\,Hz; DUTY $\in[0.25,0.75]$; $A\in[0,0.9]$\,rad; $q_0\in[-0.5,0.5]$\,rad \\
    CMA-ES & population 16, $\sigma_0=0.25$, 2500 evaluations ($10^4$ for free phase), non-zero seeds \\
    Limits & heading/roll/pitch RMS $10/10/\SI{15}{\degree}$ or $10/5/\SI{5}{\degree}$; surfacing $\le5$\,\% where stated; lift thrust share $s_L\le0.1$ for the drag-based strokes of \cref{sec:mj-liftshare} \\
    \bottomrule
  \end{tabularx}
\end{table}

\section{Model Errors Found and Corrected}

Several errors in the simulation set-up were found during the work and corrected before the reported results were produced.
They are listed because they are easy to repeat.
\begin{itemize}
  \item MJCF reads joint ranges in degrees by default; a range of $\pm1.3$ intended in radians locked every joint to
        $\pm\SI{1.3}{\degree}$. All results before this correction were discarded.
  \item Changing body masses at run time without recomputing MuJoCo's derived constants made the added mass ineffective for a
        single free body and caused a 36\,\% error in a momentum-conservation test.
  \item Visual and collision geometries imported from a URDF were both given hydrodynamic forces, doubling buoyancy, drag and added
        mass.
  \item An additive weighting of speed and attitude produced strongly rolling and yawing optima once the legs could move freely;
        it was replaced by the feasibility-first formulation \eqref{eq:feasibility}.
  \item A first measure of the lift share, $|I_L|/(|I_L|+|I_\text{drag}|)$, is about 0.5 for every gait with lift, because in steady
        swimming the lift and drag impulses cancel. Used as a limit, it could only be met by gaits that barely moved. It was
        replaced by $s_L$ of \eqref{eq:liftshare}, and the runs with the old measure were discarded.
  \item In the fluke CFD, a coarse background mesh that had not been refined around the fluke led to a failed hole cutting in which
        almost the whole background was removed; the number of hole cells is now checked automatically.
\end{itemize}
